\title{LongRoPE: Extending LLM Context Window Beyond 2 Million Tokens}

\begin{abstract}
Expanding the context window size is advantageous for large language models (LLMs). Yet, contemporary models are constrained to roughly 128k tokens, primarily because fine-tuning is expensive, long documents are not readily available, and incorporating new token positions often leads to problematic values.
\end{abstract}

This paper presents LongRoPE, which for the first time expands the context window of pre-trained LLMs to an impressive \textbf{2048k} tokens, requiring only up to 1k fine-tuning steps at training sequence lengths of 256k. Crucially, this extension preserves the original model's performance on shorter sequences. The method is based on three principal innovations: (i) we discover and utilize two distinct kinds of non-uniformities in positional interpolation through an efficient search process, enabling superior initialization during fine-tuning and an 8$\times$ extension even without additional fine-tuning; (ii) we propose a staged extension approach, wherein the LLM is first fine-tuned on 256k token sequences, after which a second positional interpolation is applied to the extended model to reach a 2048k token context window; (iii) we recalibrate LongRoPE at a sequence length of 8k to ensure high performance for short contexts. Comprehensive evaluations on LLaMA2 and Mistral across diverse benchmarks validate the advantages of our method. Models augmented with LongRoPE require only minor changes to positional embeddings and largely preserve their existing architectural optimizations. The implementation will be released at \texttt{https://github.com/microsoft/LongRoPE}
\end{abstract}

\section{Introduction}

Large Language Models (LLMs) have achieved significant advancements across numerous applications~\cite{gpt4,llama2}; however, they are constrained by a finite context window—such as the 4096 token boundary in LLaMA2~\cite{llama2}. When input length extends beyond this window, LLMs exhibit diminished performance, largely because they encounter positions outside their training distribution. Such limitations become particularly problematic for crucial use cases, including in-context learning with large numbers of examples~\cite{Huang2023FewerIM} and the operation of LLM agents~\cite{park2023generative,madaan2023selfrefine}.

Recent studies have demonstrated that the context window of pre-trained LLMs can be expanded up to approximately 128k tokens through fine-tuning with longer input sequences~\cite{longlora,pi,yarn,zhang2024soaring,liu2023scaling}. Extending the context window further faces three primary challenges. Firstly, the addition of untrained position indices can result in significant catastrophic values, thereby causing distribution shifts and impeding the convergence of the fine-tuning process~\cite{pi}. This issue is exacerbated when increasing the context from 4k to over 1000k, which incorporates more than 90\% previously unseen positions. Secondly, successful fine-tuning typically depends on availability of texts matching the desired length, but there are few datasets containing sequences longer than 1000k tokens. Furthermore, training with such extensive sequences demands immense computational effort, with substantial training duration and GPU resources required. Lastly, as the context window becomes exceedingly large, the attention mechanism becomes less effective, with tokens distributed sparsely across the sequence, ultimately leading to diminished performance on tasks with short context windows~\cite{pi}.

To address the initial challenge, a commonly employed strategy involves interpolating RoPE positional embeddings~\cite{rope,pi}, wherein position indices are rescaled into the pretrained scope, as depicted in Fig.\ref{fig:overview}. Position Interpolation (PI)~\cite{pi} adapts RoPE’s rotary angles by applying a linear interpolation based on the degree of extension. NTK~\cite{ntk,dynamicntk} proposes the use of non-uniform interpolation and extrapolation across different RoPE dimensions. YaRN~\cite{yarn} further divides RoPE dimensions into three groups according to their frequency properties, and implements extrapolation, NTK, or linear interpolation within each group. Nonetheless, the Transformer’s positional embeddings possess an intricate, non-uniform distribution of information entropy. Current methods do not fully capitalize on this subtle non-uniformity, resulting in some loss of positional information and subsequently constraining the effective context window length.

Section~\ref{sec:analysis} provides empirical evidence for two main insights: \textbf{(1)} Successful positional interpolation must account for two kinds of non-uniformities: differences across RoPE dimensions and token positions. Less interpolation is advantageous for dimensions with lower RoPE indices as well as for early token positions; however, the precise optimal interpolation strategies depend on the specific sequence length to be extended.
\textbf{(2)} Incorporating these non-uniform characteristics into the interpolation process enables preservation of essential information within the original RoPE, specifically maintaining critical dimensions and positions. This approach mitigates interpolation-related information loss, yielding improved initialization prior to fine-tuning. Additionally, it makes it feasible to achieve an 8$\times$ sequence extension even without subsequent fine-tuning.

Building on these insights, we present \textbf{LongRoPE}, a robust technique designed to expand the LLM context window past 2 \emph{million} tokens. The approach relies on three primary innovations. Firstly, {LongRoPE} capitalizes on multidimensional asymmetries in positional interpolation, determining optimized rescale factors for RoPE's rotation angles across each RoPE dimension by considering token locations. Because the parameter space for rescale factor identification grows exponentially with higher extension ratios, {LongRoPE} applies an evolutionary search process enhanced by two optimization strategies, significantly increasing search performance. Fig.~\ref{fig:overview} illustrates a sample outcome of the optimized, rescaled RoPE.

Subsequently, {LongRoPE} employs a resource-efficient, stepwise extension methodology to realize a 2048k context window, circumventing the necessity to fine-tune directly on ultra-long text sequences, which are both rare and largely inaccessible. This approach initiates by identifying a feasible 256k context length on the pre-trained LLM and carrying out fine-tuning at this scale. Leveraging our non-uniform positional interpolation, which makes it possible to expand context length by a factor of 8$\times$ without further fine-tuning, we next execute a secondary search for appropriate RoPE rescale parameters on the LLM that has already been fine-tuned for the extended context. This process ultimately secures a 2048k context window for both LLaMA2 and Mistral~\cite{mistral}.

Lastly, to address potential losses in performance when dealing with the original (shorter) context length, {LongRoPE} further fine-tunes the RoPE rescaling parameters on the expanded LLM. As with the upscaling procedure from 256k to 2048k tokens, our search algorithm is employed to rescale the 256k fine-tuned LLM for reduced context lengths of 4k and 8k, aiming to minimize positional interpolation. At inference time, if the input sequence does not exceed 8k tokens, the model applies the searched RoPE rescale factors accordingly.

A comprehensive set of experiments conducted on multiple LLMs and a range of long-context tasks highlight the efficacy of our approach. Our results indicate that LongRoPE consistently sustains low perplexity across evaluation lengths ranging from 4k to 2048k, exceeds 90\% passkey retrieval accuracy, and achieves accuracy levels on standard benchmarks within the 4096 context window that are comparable to baseline models. Additionally, LongRoPE is compatible with any LLM utilizing RoPE embeddings. We plan to make our code and LongRoPE-2048k models publicly available.

\section{Non-uniformity in Positional Interpolation}
\label{sec:analysis}

\subsection{Preliminary}

Transformer models require explicit positional information, often in the form of position embedding, to represent the order of input tokens. Our work focuses on the RoPE~\cite{rope} position embedding, which is widely used in recent  LLMs. For a token at position index $n$, its corresponding RoPE encoding can be simplified as follows:

\begin{equation}
		\fontsize{7.8}{7.8} \selectfont
	\label{eq:rope}
	\left[ cos(n\theta_0), sin(n\theta_0), cos(n\theta_1), \cdot\cdot\cdot, cos(n\theta_{d/2-1}), sin(n\theta_{d/2-1}) \right]   
\end{equation}

Here, $d$ denotes the embedding size, 
$n\theta_i$ corresponds to the rotary phase applied to the token at the $n$-th position, 
and $\theta_i = \theta^{-2i/d}$ specifies the set of rotational frequencies. RoPE typically uses a default base $\theta$ value of 10000.

\textbf{\emph{Context window extension ratio $s$} and positional interpolation}. The parameter $s$ is defined as the proportion between the new, extended context length $L'$ and the initial context length $L$, given by $s=\frac{L'}{L}$.

To extend the context window from $L$ to $L'$, current positional interpolation methods suggest downscaling rotation frequencies $\theta_i$ based on extension ratio $s$. Let $\beta=\theta^{2/d}$, and $\lambda$ denote the actual rescale factor related to $s$, we   unify these positional interpolation methods as follows:

\begin{equation}	
	\fontsize{7.5}{7.5} \selectfont
	\label{eq:unified_rope}
	\begin{aligned}
		\left[cos\left(\frac{n}{\lambda(\beta)^0}\right), sin \left(\frac{n}{\lambda(\beta)^0}\right), cos\left(\frac{n}{\lambda(\beta)^1}\right), \cdot\cdot\cdot,  sin \left(\frac{n}{\lambda(\beta)^{d/2-1}}\right) \right]   
	\end{aligned}
\end{equation}

\textbf{Linear positional interpolation (PI)}. The PI approach~\cite{pi} implements a linear scaling of position indices, constrained within the original pre-trained sequence length. Specifically, for an extension ratio $s$, the rotation angles for each position are uniformly compressed by a factor $\lambda=s$ over every RoPE dimension. This uniform contraction causes positional representations to become densely packed, which reduces the model’s effectiveness in discriminating among tokens with nearby indices. As a consequence, PI’s performance generally declines when applied at larger extension ratios.

\textbf{NTK-based interpolation and extrapolation}. Studies such as ~\cite{ntk,dynamicntk} analyze RoPE from the standpoint of information encoding, utilizing Neural Tangent Kernel (NTK) theory~\cite{jacot2018neural,tancik2020fourier}. To address the challenge of crowded positions inherent in PI, these approaches propose allocating the interpolation demands more evenly across the various RoPE dimensions. This is accomplished by scaling dimensions with higher frequencies (lower indices) less and dimensions with lower frequencies (higher indices) more, which facilitates not only positional interpolation but also extrapolation, as expressed by $\lambda = s^{i}$. The dynamic NTK model presented in~\cite{dynamicntk} further introduces adaptive scaling of the extension ratio at each position, conditioned on the ongoing sequence length. In contrast to PI, which requires model fine-tuning, NTK-based techniques are capable of extending the context window even without fine-tuning, although the extension ratio typically does not exceed 4$\times$.

\textbf{YaRN}~\cite{yarn} proposes a major advancement in the effectiveness of positional interpolation. The method segments RoPE dimensions into three distinct groups according to their frequency characteristics, applying unique interpolation techniques to each group. For the high-frequency dimensions, extrapolation is utilized ($\lambda$=1), whereas the low-frequency dimensions are interpolated linearly (PI). Dimensions that occupy the middle frequency range are treated using the NTK. A central feature of YaRN is its method for partitioning RoPE dimensions, which is currently based on manual, empirical analysis. Consequently, this approach may yield less-than-optimal results when adapting to novel LLM architectures.

\subsection{Study on Non-uniform Positional Interpolation}

Drawing upon NTK and YaRN, we observe that their performance improvements stem from the use of non-linear mechanisms, particularly the treatment of diverse frequencies within RoPE dimensions that allow for tailored interpolation and extrapolation. Nonetheless, existing approaches to non-linearity are largely predicated on heuristics crafted by researchers. This observation prompts two key inquiries: (1) Does the prevailing method for positional interpolation represent the optimal solution? (2) Are there potential non-linear functions yet to be investigated?

In order to address these questions, we employ evolution search (refer to Sec.~\ref{sec:method}) to identify improved non-uniform positional interpolation strategies for LLaMA2-7B. The process is steered by perplexity measurements, utilizing 5 randomly selected instances from the PG19~\cite{pg19} validation dataset. Our experimental investigation uncovers several notable observations.

\textbf{Finding 1}: \textit{RoPE dimensions exhibit substantial non-uniformities, which are not effectively handled by current positional interpolation methods.}

We conduct a search for the optimal value of $\lambda$ for every RoPE dimension as defined in Eq.~\ref{eq:unified_rope}. Table~\ref{tbl:exp1} presents perplexity results for LLaMA2-7B employing various methods, evaluated on the PG19 and Proof-pile~\cite{proof-pile} test datasets, without the application of fine-tuning. The solution found via our search delivers notable gains, highlighting that commonly used linear (PI) and non-uniform interpolation approaches (Dynamic-NTK and YaRN) are not optimal. Of particular interest is the observation that YaRN performs worse than both PI and NTK on PG19, as it fails to cover the intended context window length in models that have not undergone fine-tuning. Specifically, perplexity for YaRN surges after a context length of 7k tokens when the window is set to 8k.

In our investigation, we observed that the rescaling coefficients $\lambda$ in Eq.~\ref{eq:unified_rope} are distributed non-uniformly, contrasting with the constant scaling factor $s$ applied in PI, the calculation rules in NTK, and the group-based approach in YaRN. Utilizing these non-uniform scaling factors leads to substantial enhancements in LLaMA2's language modeling abilities—measured by perplexity—across context lengths of 8k and 16k, all achieved without any need for additional fine-tuning. This improvement is attributed to the modified positional embeddings, which closely retain the structure of the original RoPE, particularly preserving critical dimensions; as a result, the LLM faces less challenge in differentiating between tokens that appear in similar positions.

\textbf{Finding 2}: \textit{RoPE for the initial tokens in the input sequence should be extrapolated with less interpolation.} 

For the first $\hat{n}$ tokens of the input, we propose that reduced interpolation within RoPE is beneficial. This stems from their elevated attention scores, indicating their essential role in the attention mechanism, a trend identified in Streaming LLM~\cite{streamingllm} and LM-Infinite~\cite{han2023lminfinite}. To test this theory, we increase the context length to 8k and 16k using both PI and NTK, systematically omitting interpolation for the leading $\hat n$ tokens ($0, 2, \ldots, 256$). Setting $\hat n = 0$ yields the default PI and NTK implementations. Table~\ref{tbl:tokenposition} presents two significant findings: \textbf{(1)} preserving initial tokens without position interpolation results in measurable performance gains, and \textbf{(2)} the optimal selection for $\hat n$ varies according to the intended context window expansion.

\textbf{Finding 3}: \textit{Non-uniform positional interpolation effectively extends LLM context window in both fine-tuning and non-fine-tuning settings.} 

Our results demonstrate that applying our optimized non-uniform position interpolation yields considerable gains in extension capability for 8k and 16k contexts, even without fine-tuning. However, extending beyond these lengths necessitates fine-tuning. Consequently, we perform fine-tuning on LLaMA2-7B with our derived RoPE configuration for a 64k token context window (see Appendix for detailed experimental settings). According to Table~\ref{tbl:finetune}, our approach not only substantially surpasses PI and YaRN when used out-of-the-box with LLaMA2-7B, but also after fine-tuning. This performance advantage can be attributed to the strategic adoption of non-uniform positional interpolation, which reduces the loss of positional information and enables a more advantageous starting point for the fine-tuning process.

\textbf{Summary}. The investigation highlights two types of non-uniformity: variable RoPE dimensions and differing token positions. Leveraging these forms of non-uniformity in positional interpolation significantly boosts the context window extension performance of LLMs.

\section{LongRoPE}

\label{sec:method}
Building on these observations, we propose LongRoPE, which begins by presenting a search algorithm designed to leverage both types of non-uniformity. Utilizing this approach, we successfully scale the context window for LLMs to exceed 2 million tokens.

\subsection{Problem Formulation}

These two non-uniformities result in an extensive solution space and add layers of difficulty to the optimization process. To mitigate these challenges, we reformulate the multidimensional non-uniform position interpolation optimization challenge as a search problem.

For a LLM targeting a  context window size of $L'$ and lengthy input documents $\mathbf{X}$, where each $\mathbf{x}\in\mathbf{X}$ surpasses $L'$ in token length, we denote the original rotary angle of the $i^{th}$ dimension in RoPE embedding at token position $n$ as  $\frac{n}{\beta^i}$. The optimization problem is then formulated as follows:

\begin{equation}
\begin{gathered}
		\label{eq:problem}

\underset {\mathbf{x}\in\mathbf{X}; \, |\mathbf{x}|\ge L'}{ \operatorname {arg\,min} } \,  \mathcal{L}\, \left( \text{LLM} (\text{RoPE}, \mathbf{X})\right), \text{where } 
\\
\scalebox{0.93}{
$\underset {\begin{subarray}{l} i=0,\cdot\cdot,\frac{d}{2}-1;\\ n\in[ 0,|\mathbf{x}|);\end{subarray}}{\text{RoPE}_(n)}= {\left[\cdots, \cos\left(\mathbb{I}(\hat{\lambda}_{i}, \hat{n}) \frac{n}{\beta^i}\right), \sin\left(\mathbb{I}(\hat{\lambda}_{i}, \hat{n}) \frac{n}{\beta^i}\right),\cdots\right] }$}\\
	\text{with } \mathbb{I}(\hat{\lambda}_{i},\hat{n})=\begin{cases}
	1 &\text{if}\;  n<\hat{n} \\
	\frac{1}{\lambda_{i}} &\text{if}\; n\geq\hat{n}
	\end{cases}
	\end{gathered}
\end{equation}

In this context, we define a collection of rescale factors $\mathbb{I}(\hat{\lambda}_{i},\hat{n})$ to account for the two distinct types of non-uniformities. The variables $\hat{\lambda_i}$ and $\hat{n}$ correspond to the irregularities present in RoPE dimensions and token placement, respectively. More precisely, $\mathbb{I}(\hat{\lambda}_{i},\hat{n})$ functions to adjust the rotation angle associated with the $i^{th}$ RoPE dimension; $\hat\lambda_{i}$ serves as the scaling coefficient, while $\hat{n}$ acts as the threshold for token positions. For tokens at positions before $\hat{n}$, specifically the first $\hat{n}-1$ tokens, the factor $\hat\lambda_{i}$ is inactive, and the original rotary angle formula $\frac{n}{\beta^i}$ remains unchanged. Once the token positions reach $n \geq \hat{n}$, the scaling factor is then incorporated.

With a predetermined target context window length of $L'$, the task involves determining the set of optimal rescale factors ($\mathbb{I}(\hat{\lambda}_{0},\hat{n})$, $\mathbb{I}(\hat{\lambda}_{1},\hat{n})$, ..., $\mathbb{I}(\hat{\lambda}_{i},\hat{n})$, ...) corresponding to each dimension from the first up to the $d^{th}$ in RoPE. Through this approach, the adapted $\text{LLM}$, leveraging the rescaled $\text{RoPE}$, is able to minimize the next-token prediction loss $\mathcal{L}$ (equivalent to perplexity) over a set of input sequences $\mathbf{X}$ whose token count is $L'$.

\subsection{Searching the Non-uniform Position Interpolation}

In addressing the challenge posed by Eq.~\ref{eq:problem}, we present a straightforward and powerful approach that identifies the best RoPE rescale factors, thereby harnessing the complex multidimensional variations present in position embedding.

\textbf{Search space}. We construct an expansive search space to accommodate both forms of non-uniformity. The details of this search space are outlined in Table~\ref{tbl:searchspace}. In particular, our approach permits independent selection of a rescale factor for every dimension within the RoPE embedding. To streamline the configuration of the search space, we focus on searching for $\lambda_i$ and $\hat{n}$, rather than the composite variable $\mathbb{I}(\hat{\lambda}_{i},\hat{n})$, where  $\hat{\lambda}_{i}=1/\lambda_i$. 
As indicated in Table~\ref{tbl:searchspace}, the parameter $\lambda_i$ can take values ranging from 1.0 (which corresponds to straightforward extrapolation) up to $s\times1.25$ (enabling broader interpolation than PI), incremented by 0.01 for each search trial. Here, $s$ denotes the proportion by which the context window is extended.

The parameter $\hat{n}$ specifies how many of the earliest token positions maintain the unaltered RoPE embedding, foregoing any position interpolation. In our experiments, $\hat{n}$ is selected from the set \{0, 1, 2, 4, 8, 12, 16, 20, 24, 28, 32, 64, 128, 256\}. Setting $\hat{n}=0$ implies that every token position utilizes the optimized rescale factors.

\textbf{Evolution-based search}. The search space outlined in Table~\ref{tbl:searchspace} encompasses a wide array of positional interpolation strategies, which presents substantial difficulties in achieving efficient search. As an illustration, for a $s=4\times$ scaling, the number of possible configurations amounts to $400^{128/2}\times14=4\times10^{167}$. This exponential increase in search space with higher extension ratios intensifies the complexity. To mitigate this, we employ an evolution search algorithm~\cite{guo2020single}, supplemented by two optimization methods that considerably enhance the effectiveness of our search process. The general workflow of the search is provided in Algorithm~\ref{alg:search}.

\textit{Optimized initial population generation}. In place of purely random initialization for the $P$ rescale factors, we explicitly include the three RoPE rescale factors that represent PI, NTK, and YaRN in the starting population as distinct individuals. The rest of the individuals, totaling $P$-3, are generated by applying random mutations to these three rescale factors with a probability $p$.

\textit{Monotonically non-decreasing constraint}. Once the initial population has been generated, we evaluate each individual by calculating the LLM perplexity. For this, the respective RoPE rescale factors are assigned to the target LLM, and the perplexity over input $\mathbf{X}$ is measured. The $k$ highest-scoring individuals are then selected as parents to guide the evolutionary process. Nonetheless, the extensive search space means that random mutation and crossover can frequently produce suboptimal candidates, resulting in a surplus of perplexity evaluations that offer little benefit. This inefficiency becomes more pronounced as $L'$ increases, due to the computationally demanding nature of each perplexity assessment.

To mitigate this issue, we enforce a monotonicity constraint on the RoPE rescaling factors such that $\lambda_i\leq \lambda_{i+1}$ for all $i$. Only those RoPE configurations that conform to this criterion are evaluated for LLM perplexity, thereby streamlining and reducing the computational expense associated with the search. In particular, this means that $\lambda_i$ must increase with respect to the RoPE dimension index ($i=0, \ldots, 63$). The rationale for imposing monotonicity across dimensions is inspired by NTK theory~\cite{jacot2018neural,tancik2020fourier,ntk}, which indicates that dimensions encoding higher frequencies (lower indices) require less interpolation (corresponding to a smaller $\lambda_i$), while dimensions associated with lower frequencies (higher indices) benefit from greater interpolation capability (implying a larger $\lambda_i$).

\begin{algorithm}[t]
	\small
	\caption{The search algorithm}
	\textbf{Input:} target LLM, input samples $\mathbf{X}$,   population size $P$, mutation size $N_1$, crossover size $N_2$, max iterations $\mathcal{T}$, mutate probability $p$\\

	\begin{algorithmic}[1]
		\label{alg:search}
		\STATE Top-k$\leftarrow \phi$;	
		\STATE $\text{P}_0 \leftarrow$ \textit{Initialize\_population\_with\_optimization} ($P$, $\mathbf{X}$, $p$); 
		\FOR{$i = 1$ to $\mathcal{T}$}
		\STATE \textit{Compute\_perplexity} (LLM, $\text{P}_{i-1}$, $\mathbf{X}$);
		\STATE  Top-k $\leftarrow$ \textit{Update\_Topk} (Top-k, $\text{P}_{i-1}$);
		\STATE$\text{P}_{mutation} \leftarrow$ \textit{Mutation\_with\_mono\_constraint} (Top-k, $N_1$, $p$); 
		\STATE $\text{P}_{crossover} \leftarrow$ \textit{Crossover\_with\_mono\_constraint} (Top-k, $N_2$);
		\STATE $\text{P}_i \leftarrow \text{P}_{mutation} \cup \text{P}_{crossover} \cup$ Top-k;
		\ENDFOR
		\STATE Output the member from Top-k with minimal perplexity;
	\end{algorithmic}
\end{algorithm}

\textbf{8$\times$ expansion absent fine-tuning}. Utilizing an evolutionary search process, our approach reliably discovers optimal non-uniform RoPE scaling factors, retaining important dimensions and positional information so as to reduce loss from interpolation. As shown in Fig.\ref{fig:non-ft}, this technique enables LLaMA2’s context window to be increased from 4k up to 32k without requiring any fine-tuning. Meanwhile, prior approaches—including PI, as well as non-uniform NTK and YaRN—demonstrate a sharp rise in perplexity after doubling the window size.

\subsection{Extending LLM Context Window to 2048K}
\label{sec:extendto2048k}

\textbf{Progressive extension to 2048k}. In this section, we present our approach for expanding the context window of pre-trained LLMs beyond the conventional 4k, reaching lengths exceeding 2048k. Our method of non-uniform positional interpolation enables an 8$\times$ context length expansion without the need for additional fine-tuning. However, when seeking even greater extensions (for example, up to 512$\times$), fine-tuning becomes indispensable. A potential strategy is to identify the appropriate RoPE rescaling factors corresponding to the 2048k target window size, followed by fine-tuning the model. Despite this, substantial obstacles remain, chiefly the considerable computational cost associated with such large-scale training. Furthermore, our practical attempts indicate that fine-tuning LLMs for extremely large extension ratios proves difficult to execute effectively (refer to Appendix for further details).

Fortunately, LongRoPE demonstrates efficacy when applied to both baseline and fine-tuned variants of the extended LLM. To this end, we propose a streamlined, incremental approach that reaches the desired 2048k context window in only 1k fine-tuning iterations, utilizing a training length capped at 256k.

$\Diamond$ \textit{Expanding pre-trained LLM to 256k tokens utilizing LongRoPE search}. Using LLaMA2 for illustration, we perform a search process targeting context window lengths of 128k and 256k. The corresponding extension ratios at these intervals are 32$\times$ and 64$\times$, in order.

$\Diamond$ \textit{Fine-tuning to 256k}. In the subsequent phase, we adapt the pre-trained LLM to extend its context window to 256k. This process begins with fine-tuning LLaMA2 over 400 steps utilizing RoPE rescaling factors aligned with a 128k context. After this initial stage, we update the checkpoint to incorporate RoPE rescaling factors for 256k, followed by a further 600 fine-tuning steps. Compared to directly fine-tuning for 256k, this staged approach offers improved efficiency.

$\Diamond$ \textit{Expanding fine-tuned extended LLMs to a 2048k context window via LongRoPE search}. In the last step, a subsequent search is conducted using the LLM that has previously been fine-tuned for a context length of 256k. This process yields a remarkably large context window of 2048k, achieved without additional fine-tuning. The ultimate window extension reaches a $512\times$ increase.

\textbf{Restoration of performance in shorter context windows}. Following the expansion to a substantial 2048k context window, a decline in performance is observed in the initial context window. This phenomenon is well-documented in the literature on positional interpolation~\cite{pi}, which compresses position embeddings in the lower region of the original context window, resulting in less effective use of embedding space and degraded model performance. Specifically, with an extension ratio of 512$\times$, the positional embeddings for the original 4k window are densely packed, further exacerbating this issue.

To address this issue, we carry out an additional evolutionary search on the expanded LLM, tuning the RoPE rescale parameters specifically for shorter context lengths such as 4k and 8k. The upper limit for searched $\lambda$ is lowered, since less positional interpolation is necessary when dealing with briefer contexts. While performing inference, the LLM modulates the relevant RoPE rescale factors in real time.

\section{LongRoPE}

\label{sec:method}
Building on these observations, we propose LongRoPE, which incorporates an optimized search algorithm designed to leverage the identified non-uniformities. Utilizing this approach, we successfully expand the LLM context window to exceed 2 million tokens.

\subsection{Problem Formulation}

The presence of these two non-uniformities significantly expands the solution space and complicates the optimization process. To manage these challenges, we recast the multidimensional non-uniform position interpolation optimization task into a search-based problem.

For a LLM targeting a  context window size of $L'$ and lengthy input documents $\mathbf{X}$, where each $\mathbf{x}\in\mathbf{X}$ surpasses $L'$ in token length, we denote the original rotary angle of the $i^{th}$ dimension in RoPE embedding at token position $n$ as  $\frac{n}{\beta^i}$. The optimization problem is then formulated as follows:

\begin{equation}
\begin{gathered}
		\label{eq:problem}

\underset {\mathbf{x}\in\mathbf{X};\, |\mathbf{x}|\ge L'}{\operatorname{arg\,min}}\; \mathcal{L}\left(\text{LLM}(\text{RoPE},\,\mathbf{X})\right), \text{where}
\\
\scalebox{0.93}{
$\underset{\begin{subarray}{l}i=0,\,\cdots,\,\frac{d}{2}-1;\\ n\in[0,|\mathbf{x}|);\end{subarray}}{\text{RoPE}_\,(n)} = \left[\ldots, \cos\left(\mathbb{I}(\hat{\lambda}_{i},\,\hat{n})\,\frac{n}{\beta^i}\right),\,\sin\left(\mathbb{I}(\hat{\lambda}_{i},\,\hat{n})\,\frac{n}{\beta^i}\right),\ldots\right]$
}
\\
\text{where}\; \mathbb{I}(\hat{\lambda}_{i},\,\hat{n}) = \begin{cases}
	1 & \text{if}\; n<\hat{n} \\
	\frac{1}{\lambda_{i}} & \text{if}\; n \geq \hat{n}
\end{cases}
	\end{gathered}
\end{equation}

In this approach, we define a set of rescaling functions, $\mathbb{I}(\hat{\lambda}_{i},\hat{n})$, designed to account for two distinct types of non-uniformity. The parameter $\hat{\lambda_i}$ corresponds to non-uniformity in the RoPE dimensions, while $\hat{n}$ captures irregularities associated with token positions. More concretely, $\mathbb{I}(\hat{\lambda}_{i},\hat{n})$ is utilized to adjust the rotation angle of the $i^{th}$ RoPE dimension, wherein $\hat\lambda_{i}$ serves as the scaling factor and $\hat{n}$ sets the threshold for token positions. For token indices less than $\hat{n}$, the scaling factor $\hat\lambda_{i}$ is not applied; thus, the conventional RoPE rotary angle $\frac{n}{\beta^i}$ is retained. Once token positions reach or surpass $n \geq \hat{n}$, the rescale factor comes into effect.

With a predetermined target context window length of $L'$, the task is to determine the set of optimal rescaling factors ($\mathbb{I}(\hat{\lambda}_{0},\hat{n})$, $\mathbb{I}(\hat{\lambda}_{1},\hat{n})$, ..., $\mathbb{I}(\hat{\lambda}_{i},\hat{n})$, ...) for each RoPE dimension from the $1^{st}$ up to the $d^{th}$. These rescaled values are intended to modify the target $\text{LLM}$’s $\text{RoPE}$ representations such that the model obtains the lowest possible next-token prediction error, denoted $\mathcal{L}$ (or perplexity), when processing input sequences $\mathbf{X}$ of length $L'$.

\subsection{Searching the Non-uniform Position Interpolation}

In order to address the challenge outlined in Eq.~\ref{eq:problem}, we present a straightforward and powerful approach that identifies the ideal RoPE rescaling parameters. This enables the method to capitalize on the complex, non-uniform characteristics present across multiple dimensions in positional embeddings.

\textbf{Search space}. We construct an extensive search space that captures both types of non-uniformity. The search space is presented in Table~\ref{tbl:searchspace}. Our approach enables tuning of a distinct rescale factor for each axis in the RoPE embedding. To simplify the search process, we optimize over $\lambda_i$ and $\hat{n}$, rather than $\mathbb{I}(\hat{\lambda}_{i},\hat{n})$, noting that $\hat{\lambda}_{i}=1/\lambda_i$. Table~\ref{tbl:searchspace} specifies that the range for $\lambda_i$ spans from 1.0 (which corresponds to direct extrapolation) up to $s\times1.25$ (representing more substantial interpolation relative to PI) with increments of 0.01, where $s$ indicates the context window extension ratio being targeted.

The parameter $\hat{n}$ specifies how many leading token positions preserve the original RoPE embedding, foregoing position interpolation. In practice, we select $\hat{n}$ from the set \{0, 1, 2, 4, 8, 12, 16, 20, 24, 28, 32, 64, 128, 256\} during the search process. If $\hat{n}=0$, it indicates that every token position applies the rescale factors identified during searching.

\textbf{Evolution-based search}. The search space outlined in Table~\ref{tbl:searchspace} incorporates a wide variety of positional interpolation approaches, creating a substantial obstacle for systematic exploration. For instance, increasing the extension ratio to $s=4\times$ results in $400^{128/2}\times14$=4$\times10^{167}$ possible configurations. This exponential growth in search space is exacerbated as the extension ratio rises. To overcome these challenges, we employ evolution search~\cite{guo2020single}, supplemented by two additional optimization strategies that significantly enhance search effectiveness. The complete search workflow is depicted in Algorithm~\ref{alg:search}.

\textit{Enhanced initial population formation}. Rather than randomly assigning values to all $P$ rescale factors at the start, we explicitly incorporate the three RoPE rescale factors associated with PI, NTK, and YaRN as members of the initial population. The other $P$-3 individuals are produced by introducing random mutations to these three rescale factors with a mutation probability of $p$.

\textit{Monotonically non-decreasing constraint}. Once the initial population is created, the LLM perplexity is evaluated for every individual by implementing the associated RoPE rescale factors in the target LLM and measuring the perplexity on input $\mathbf{X}$. The highest-ranked $k$ individuals are then selected as parents for the evolutionary process. Nevertheless, due to the extensive search space, simple mutation and crossover methods can frequently sample suboptimal candidates, incurring excessive and redundant perplexity evaluations. This inefficiency becomes more pronounced with larger $L'$, as each perplexity computation involves costly inference steps.

To mitigate this issue, we enforce a non-decreasing monotonicity condition on the RoPE rescaled factors, such that $\lambda_i \leq \lambda_{i+1}$. Only those RoPE configurations adhering to this criterion are considered for perplexity assessment in the LLM, which streamlines the search process considerably. Specifically, the sequence of $\lambda_i$ values is constrained to be monotonically increasing across RoPE dimensions, with $i = 0, \ldots, 63$. This constraint leverages insights from NTK theory~\cite{jacot2018neural,tancik2020fourier,ntk}: lower-index dimensions correspond to higher frequencies that necessitate minimal interpolation (implying smaller $\lambda_i$), while higher-index dimensions align with lower frequencies, thus allowing more interpolation (reflected in larger $\lambda_i$).

\begin{algorithm}[t]
	\small
	\caption{The search algorithm}
	\textbf{Input:} target LLM, input samples $\mathbf{X}$, population size $P$, mutation size $N_1$, crossover size $N_2$, maximum iterations $\mathcal{T}$, mutation probability $p$ \\

	\begin{algorithmic}[1]
		\label{alg:search}
		\STATE Top-k $\leftarrow$ $\phi$;
		\STATE $\text{P}_0$ $\leftarrow$ \textit{Initialize\_population\_with\_optimization}($P$, $\mathbf{X}$, $p$);
		\FOR{$i = 1$ to $\mathcal{T}$}
		\STATE \textit{Compute\_perplexity}(LLM, $\text{P}_{i-1}$, $\mathbf{X}$);
		\STATE Top-k $\leftarrow$ \textit{Update\_Topk}(Top-k, $\text{P}_{i-1}$);
		\STATE $\text{P}_{mutation}$ $\leftarrow$ \textit{Mutation\_with\_mono\_constraint}(Top-k, $N_1$, $p$);
		\STATE $\text{P}_{crossover}$ $\leftarrow$ \textit{Crossover\_with\_mono\_constraint}(Top-k, $N_2$);
		\STATE $\text{P}_i$ $\leftarrow$ $\text{P}_{mutation}$ $\cup$ $\text{P}_{crossover}$ $\cup$ Top-k;
		\ENDFOR
		\STATE Return the candidate from Top-k with minimum perplexity;
	\end{algorithmic}
\end{algorithm}

\textbf{8$\times$ extension achieved without fine-tuning}. Utilizing evolutionary search, our approach selects non-uniform RoPE rescale factors, maintaining crucial dimensions and positions so as to reduce information degradation caused by interpolation. As illustrated in Fig.\ref{fig:non-ft}, this strategy successfully enlarges LLaMA2's context window from 4k up to 32k without the need for fine-tuning. In comparison, previously established methods, including PI, as well as non-uniform NTK and YaRN, result in a significant increase in perplexity once the extension surpasses 2$\times$.

\subsection{Extending LLM Context Window to 2048K}
\label{sec:extendto2048k}

\textbf{Progressive extension to 2048k}. Here, we present our approach for scaling the context window of pre-trained LLMs beyond the conventional 4k, reaching up to more than 2048k. Our non-uniform positional interpolation method enables an 8$\times$ expansion without requiring fine-tuning. However, achieving much larger increases (such as 512$\times$) necessitates fine-tuning. One possible strategy involves identifying suitable RoPE rescaled factors specifically for the desired 2048k length and subsequently conducting fine-tuning. This, however, is hampered by the extremely high demand on computational resources. In addition, our practical experience indicates that fine-tuning LLMs at such high extension rates presents significant difficulties (see Appendix).

LongRoPE proves to be beneficial not only for the base model but also for extended LLMs after fine-tuning. Consequently, we propose a streamlined and incremental approach that enables reaching the desired 2048k sequence length by performing only 1k fine-tuning steps at a training length of up to 256k.

$\Diamond$ \textit{LongRoPE search for expanding pre-trained LLM context to 256k}. Using LLaMA2 as a case study, we explore extension strategies to reach context window sizes of 128k and 256k. This involves scaling up the original context window by factors of 32$\times$ and 64$\times$, accordingly.

$\Diamond$ \textit{Fine-tuning to 256k}. After initial pre-training, we adapt the LLM to accommodate a context window of 256k tokens. This is accomplished by first fine-tuning LLaMA2 over 400 iterations with RoPE rescaling parameters set for 128k. Subsequently, we switch the RoPE rescaling to match 256k and carry out a further 600 fine-tuning steps on the resultant model checkpoint. This staged approach demonstrates superior efficiency compared to immediately fine-tuning the model for a 256k context window.

$\Diamond$ \textit{Applying LongRoPE search to augment a fine-tuned extended LLM to 2048k}. As a concluding step, we conduct an additional search phase on the LLM previously fine-tuned for a context length of 256k. This process enables the model to reach a substantial context window of 2048k, without necessitating extra rounds of fine-tuning. The achieved extension represents an increase of $512\times$.

\textbf{Restoring performance in shorter context windows}. Upon increasing the context window to an extensive 2048k size, we observe diminished performance within the model's initial context span. This degradation is well documented in positional interpolation research~\cite{pi}, since it compresses high-dimensional position embeddings in the original window into a far tighter space, which hampers the effectiveness of the language model. Specifically, under a 512$\times$ scaling factor, the positional encodings within the baseline 4k context window are highly congested.

To address this issue, we conduct an additional evolutionary search on the expanded LLM to fine-tune RoPE rescale parameters specifically for shorter context lengths, such as 4k and 8k. Because these lengths require less positional interpolation, we restrict the upper bound of $\lambda$ during the search process. When performing inference, the LLM modifies the relevant RoPE rescale factors in real time.

 \section{Experiments}

\label{sec:exp}
\subsection{Setup}
\textbf{Evaluation Tasks and models}. Our experiments utilize LongRoPE with LLaMA2-7B and Mistral-7B, assessing performance in three dimensions: (1) perplexity for long-context LLMs on extended documents; (2) a Passkey retrieval challenge gauging the model’s capacity to locate a straightforward passkey amid copious unrelated content; and (3) conventional LLM benchmark evaluation restricted to a 4096-token context window.

\textbf{Fine-tuning}. For LLaMA2, a learning rate of $2 \times 10^{-5}$ with linear decay is applied, and training proceeds with a global batch size of 32. The model is initially fine-tuned for 400 steps using the Redpajama~\cite{redpajama} dataset, which is divided into 128k-token segments, each surrounded by BOS and EOS tokens. Afterward, utilizing the final checkpoint from this stage, the context window is extended to 256k through an additional 600 training steps. The 128k context configuration is trained on 8 A100 GPUs with the distributed training framework~\cite{cube}, whereas the 256k context window necessitates 16 A100 GPUs. For Mistral, a fixed learning rate of $1 \times 10^{-6}$ is adopted along with a global batch size of 64. The protocol outlined in YaRN~\cite{yarn} is followed for both the 128k and 256k models, involving 400 steps using Together Computer’s Long-Data Collections~\cite{mistraldata} with sequences of 16k length. Training for Mistral is carried out on 4 A100 GPUs.

\textbf{Search}. When the target window size is no greater than 256k, we set $P$=64, $N_1$=$N_2$=16, $p$=0.3, and $\mathcal{T}$=40, choosing the top 32 candidates for mutation and crossover during each iteration. To evaluate perplexity, 5 random samples are drawn from the PG19 validation set, each meeting at least the target context length. For target windows exceeding 512k, the number of individuals in the population, as well as the sizes for mutation and crossover, are reduced by half. In this regime, perplexity is assessed using 3 randomly chosen samples from the Pile-Books3~\cite{pile} validation data.

\textbf{Baselines}. To achieve a context window of 2048k, we applied fine-tuning to models with 128k and 256k windows. Accordingly, LongRoPE-2048k (ft=128k) was produced for LLaMA2, and LongRoPE-2048k (ft=256k) for Mistral. The quartet of models is evaluated alongside prominent baselines in context window expansion, specifically including publicly available LLMs trained with positional interpolation through PI, NTK, and YaRN approaches. The set of comparisons covers Together-32k~\cite{together}, Code LLaMA~\cite{codellama}, LongLoRA-full-FT-100k~\cite{longlora}, as well as YaRN-LLaMA and YaRN-Mistral~\cite{yarn}.

\subsection{Main Results}

\label{sec:mainresult}
\textbf{Language modeling on long sequences up to 256k tokens}. Our initial analysis focuses on benchmarking performance against leading long-context LLMs at a context window of 256k tokens. To assess generalization, we utilize the test splits from two benchmarks: Proof-pile~\cite{pg19} and PG19~\cite{pile}. Perplexity is measured at different input sequence lengths using a sliding window of size 256. For the PG19 dataset, the entire test split comprising 100 documents is evaluated. Following the protocol from YaRN~\cite{yarn}, we sample 10 instances from the Proof-pile set, ensuring that each selected sample exceeds 128k tokens in length.

Table~\ref{tbl:proof-pile} and Table~\ref{tbl:pg19} present a comparison of the perplexity scores for LLaMA2 and Mistral when extended with various interpolation approaches on the Proof-pile and PG19 benchmarks, respectively. Two main findings emerge: \textbf{(1)} our extended models consistently demonstrate lower perplexity as the evaluation length increases from 4k up to 256k, indicating improved utilization of extended context. \textbf{(2)} Despite operating in a setting where the context window is stretched 16$\times$ beyond typical limits—an especially difficult scenario for preserving short-range performance—our LongRoPE-2048k configurations achieve better results than state-of-the-art baselines for context windows up to 256k.

\textbf{Modeling long language sequences surpassing 2000k tokens}. To assess performance with exceptionally lengthy texts, we employ the Books3~\cite{pile} dataset. To streamline evaluation, we randomly choose 20 books, each with a length greater than 2048k tokens, and apply a sliding window approach using a size of 256k.

As summarized in Table~\ref{tbl:books3}, LongRoPE is able to broaden the context window for both LLaMA2-7B and Mistral-7B up to 2048k tokens, while maintaining or surpassing the perplexity scores of baseline models at shorter context spans (8k-128k). Our results further reveal marked differences in model behavior between the 2048k LLaMA2 and Mistral variants. While Mistral delivers superior performance compared to baselines on shorter contexts, its perplexity rises above 7 for sequences exceeding 256k tokens. In the case of LLaMA2, performance trends are as expected: perplexity continues to improve as the context length grows, with only minor increases observed at 1024k and 2048k tokens. Additionally, for LLaMA2, LongRoPE-2048k shows improved results when fine-tuned with a window size of 256k compared to 128k, attributable to a reduced secondary extension ratio (8$\times$ versus 16$\times$). Meanwhile, Mistral achieves better outcomes when fine-tuned at 128k. The underlying reason is related to our adherence to the YaRN configuration for Mistral, using a 16k training sequence length for both 128k and 256k fine-tuning settings, which limits Mistral's capacity for further extending its context window following fine-tuning.

\textbf{Passkey retrieval}. Next, we investigate the practical limits of the context window in generation-based tasks. For this, we adopt the synthetic passkey retrieval evaluation introduced by~\cite{passkey}. Here, the task for the model is to locate and extract a randomly generated five-digit passkey embedded within an extended document. The prompt structure can be found in the appendix. To assess performance, we conduct 10 runs of the retrieval task, ensuring that the passkey's position is randomly and uniformly distributed throughout the evaluation context length in each iteration.

Fig.~\ref{fig:passkey} illustrates the retrieval accuracy comparison against baseline models. The performance of existing approaches declines steeply to zero once the sequence length exceeds 128k tokens. Conversely, LongRoPE-LLaMA2-2048k (ft=256k) consistently achieves strong retrieval results ($\ge$90\%) across the entire range from 4k to 2048k tokens, despite the difficulty of identifying a passkey among millions of candidates. LongRoPE-Mistral-2048k (ft=128k) retains perfect accuracy up to 1800k tokens, with a reduction to 60\% at 2048k, mirroring the trend in Table~\ref{tbl:books3} where perplexity exhibits a slight uptick at the longest context.

\textbf{Evaluation on conventional benchmarks within the native context window}. LongRoPE-2048k models are assessed on their original context window, leveraging the Hugging Face Open LLM Leaderboard~\cite{huggingfaceboard} in both zero-shot and few-shot paradigms. Specifically, 25-shot ARC-Challenge~\cite{allenai:arc}, 10-shot HellaSwag~\cite{zellers2019hellaswag}, 5-shot MMLU~\cite{mmlu}, and 0-shot TruthfulQA~\cite{lin2021truthfulqa} tasks are employed.

Table~\ref{tbl:commonsense} indicates that our models perform similarly to the established benchmark, which was intended for a limited context window, and actually surpass the original Mistral on TruthfulQA by a margin of +0.5\%. Although LongRoPE-LLaMA2-2048k, which was fine-tuned at 256k, exhibits marginally greater decline in performance, it still maintains acceptable results across the majority of tasks.

\subsection{Ablation Results}

\textbf{Evaluating the impact of the second positional interpolation}. As part of our progressive extension approach, the search algorithm is leveraged to apply a second, non-uniform positional interpolation to the extended, fine-tuned LLMs. To assess its impact, we performed experiments utilizing our fine-tuned LLaMA2-256k model. The model was further extended to lengths of 512k, 1024k, and 2048k through both PI and YaRN. As indicated in Table~\ref{tbl:secondsearch}, our method using non-uniform positional interpolation is able to maintain stable perplexity values. By comparison, employing PI and YaRN results in a rapid increase in perplexity as the extension factor grows.

\textbf{Recovery at short context lengths.} In order to address degradation in performance at reduced context sizes, we re-optimize the RoPE scaling parameters for LongRoPE-2048k using our established search methodology. During this process, we impose stricter upper bounds on scale factors, intentionally minimizing interpolation when the model operates at 4k and 8k contexts. The outcomes in Table~\ref{tbl:shortperformance} present a perplexity analysis for LongRoPE-LLaMA2-2048k evaluated on Proof-pile at both 4k and 8k lengths, in addition to the mean accuracy across standard LLM benchmarks. As depicted, the proposed adjustment substantially boosts performance when working with short context sequences.

\textbf{Examination of two types of non-uniformities}. To determine the effect of each non-uniformity on performance, we conducted ablation studies. Our experimental setup includes: (i) expanding LLaMA2-7B to shorter sequence lengths of 16k and 32k using various approaches—PI, searching solely for the optimal RoPE dimension, and jointly optimizing both non-uniformities; (ii) scaling a fine-tuned LLaMA2 model with 256k sequence length to 2048k using the same methodology. Perplexity scores were measured prior to any additional fine-tuning. Results in Table~\ref{tbl:searchdimension} indicate that introducing non-uniformity within the RoPE dimension notably lowers perplexity relative to the linear interpolation used in PI. Incorporating non-uniformity in token position yields a marked benefit at 16k and 32k, whereas its effect diminishes at the 2048k level, likely attributed to the challenges posed by such extensive sequence lengths. Simply retaining original token positions without interpolation offers little utility here; exploration of this aspect remains an open avenue for subsequent research.

\section{Related Works}

Alongside techniques that employ position interpolation, this section examines other relevant methods in the literature.

\textbf{Retrieval-based} techniques rely on external memory modules to store extensive historical contexts, and utilize retrieval mechanisms to access pertinent documents during inference~\cite{longllama,wang2023augmenting,borgeaud2022improving}. Such approaches usually require significant adjustments to the internal structures of large language models (LLMs). In comparison, our approach introduces only minimal changes to positional embeddings, offering a more streamlined solution. Furthermore, it enables a wider range of applications beyond retrieval, including long-form document summarization and few-shot learning scenarios.

\textbf{Attention-based context window extensions}. In addition to methods like positional embedding interpolation, several studies have managed to extend the effective input context by modifying attention mechanisms within the standard LLM context window size~\cite{han2023lminfinite, streamingllm,parallelwindow}. Central to these approaches is the use of innovative attention masks designed to control the growth of attention computation brought on by additional positions. These approaches are complementary to positional interpolation techniques.

\textbf{Fine-tuning based approaches} center on strategies to adapt pre-trained LLMs for extended context windows by modifying position embeddings. Methods such as Code LLaMA~\cite{codellama}, LLaMA2 Long~\cite{xiong2023effective}, and ScaledRoPE~\cite{liu2023scaling} typically select a significantly enlarged base value for RoPE before conducting fine-tuning on the desired sequence length. In contrast, our approach enables adaptability across multiple target lengths and scales up to sequences exceeding 2M tokens. Recently, the computational burden of fine-tuning for very long contexts (e.g., more than 128k tokens) has prompted the development of methods like LongLoRA~\cite{longlora} and PoSE~\cite{zhu2023pose}, which aim to reduce associated GPU costs. Our method operates independently and can be combined with such efficient fine-tuning solutions.

\section{Conclusion}

This paper introduces LongRoPE, a technique that substantially increases the context length of LLMs to an extraordinary 2048k tokens, all while preserving their performance in the original, shorter context window. Our approach leverages two distinct types of non-uniformity within RoPE positional embeddings, optimized through a streamlined evolutionary search procedure. This methodology yields dual advantages: it not only offers optimal initial parameters for subsequent fine-tuning but also allows for an 8$\times$ expansion of the context window in the absence of fine-tuning. Further, we employ a gradual extension process that utilizes LLMs fine-tuned at 256k token lengths, successfully achieving a 2048k context window without requiring additional fine-tuning. Comprehensive experimental results demonstrate LongRoPE’s effectiveness. We anticipate that our LongRoPE-2048k model will unlock numerous applications demanding lengthy contexts and prompt future advancements in the field.

\section*{Broader Impacts} This study aims to contribute to ongoing developments within the Machine Learning community. While our research may carry various societal implications, we do not identify any particular issues that warrant explicit discussion at this time.

	\balance

\end{document}